This is the simplest case for a FE code: a 1D (temperature) diffusion problem. 
It puts into practice what is presented in section \ref{sec:diff1D}.
The initial temperature profile is as follows:

\begin{center}
\includegraphics[width=10cm]{python_codes/fieldstone_42/images/tempinit}
\end{center}

\[
T(x,t=0)=200 \quad x<L_x/2
\quad\quad\quad
T(x,t=0)=100 \quad x\geq L_x/2
\]

The properties of the material are as follows:
$\rho=3000$, $k=3$, $C_p=1000$ and the domain size is $L_x=100\text{km}$.
Boundary conditions are $T(t,x=0)=200\degree C$ and $T(t,x=L_x)=100\degree C$.
There are {\tt nelx} elements and {\tt nnx} nodes. 
All elements are {\tt hx} long. 
The code will carry out {\tt nstep} timesteps of length {\tt dt} or will stop before that
when steady state is reached.
The code structure is summarised hereunder:
\begin{center}
\includegraphics[width=8cm]{python_codes/fieldstone_42/images/flowchart}
\end{center}

